\documentclass[letterpaper,12pt]{article}
\usepackage[paperheight=27.94cm,paperwidth=21.59cm,bindingoffset=0in,left=3cm,right=2.0cm, top=3.5cm,bottom=2.5cm, headheight=20pt, headsep=1.0\baselineskip]{geometry}
\usepackage{graphicx,lastpage}
\usepackage{upgreek}
\usepackage{censor}
\usepackage[spanish,es-tabla]{babel}
\usepackage{pdfpages}
\usepackage{tabularx}
\usepackage{graphicx}
\usepackage{adjustbox}
\usepackage{xcolor}
\usepackage{colortbl}
\usepackage{rotating}
\usepackage{multirow}
\usepackage[utf8]{inputenc}
\usepackage{float}
\usepackage{hyperref}
\usepackage{listings}
\usepackage{xurl}

\setlength{\intextsep}{6pt}
\setlength{\abovecaptionskip}{4pt}
\setlength{\belowcaptionskip}{2pt}

\renewcommand{\tablename}{Tabla}
\usepackage{fancyhdr}
\pagestyle{fancy}
\fancyhf{}
\lhead{Criptografía y Seguridad en Redes}
\rhead{Laboratorio 2}
\cfoot{Página \thepage\ de \pageref{LastPage}}

\graphicspath{{../evidencias/capturas/}}

\lstset{
  basicstyle=\ttfamily\small,
  breaklines=true,
  columns=fullflexible,
  frame=single,
  showstringspaces=false,
  keepspaces=true
}

\newcommand{\labfigure}[4][\textwidth]{%
  \begin{figure}[H]
    \centering
    \includegraphics[width=#1,height=0.72\textheight,keepaspectratio]{#2}
    \caption{#3}
    \label{#4}
  \end{figure}
}

%
\begin{document}
%
   \title{\Huge{Informe Laboratorio 2}}

   \author{\textbf{Sección 3} \\  \\Edicson Solar Salinas \\ e-mail: edicson.solar@mail.udp.cl}
          
   \date{Septiembre de 2026}

   \maketitle
   
   \tableofcontents
 
  \newpage

\section{Descripción de actividades}

En este laboratorio trabajé con DVWA dentro de un contenedor local. El servicio quedó
publicado solamente en \texttt{127.0.0.1:6969}, por lo que todas las solicitudes de burp
suite, curl e hydra se dirigieron al mismo formulario vulnerable sin salir del equipo. El
objetivo fue ejecutar fuerza bruta de forma controlada, comprobar accesos válidos e
inválidos y después comparar el tráfico de las tres herramientas en wireshark.

Las actividades fueron las siguientes.

\begin{itemize}
  \item Desplegar \href{https://github.com/digininja/DVWA}{DVWA} mediante docker y
  explicar los parámetros utilizados.
  \item Interceptar el formulario \texttt{/vulnerabilities/brute/} con burp suite,
  configurar intruder y obtener al menos dos pares válidos.
  \item Copiar desde las herramientas de desarrollador un comando curl, ejecutar un
  acceso válido e inválido y demostrar cinco diferencias.
  \item Ejecutar hydra sobre el mismo formulario con diccionarios acotados y obtener al
  menos dos pares válidos.
  \item Capturar el tráfico de curl, burp e hydra, explicar paquetes representativos,
  comparar sus características y evaluar si el software puede reconocerse desde la red.
\end{itemize}

\section{Desarrollo de actividades según criterio de rúbrica}

\subsection{Levantamiento de docker para correr DVWA (dvwa)}

Al comenzar, el servicio de docker estaba inactivo. Lo inicié con
\texttt{sudo systemctl start docker} y comprobé que la imagen
\texttt{vulnerables/web-dvwa:latest} ya estaba disponible localmente. Por esa razón usé
\texttt{--pull=never}, la ejecución no necesitaba descargar ni sustituir la imagen
existente. El contenedor se creó con este comando:

\begin{lstlisting}
docker run --pull=never --name dvwa-lab2 --restart=no -d \
  -p 127.0.0.1:6969:80 vulnerables/web-dvwa:latest
\end{lstlisting}

\texttt{--name dvwa-lab2} asigna un nombre reconocible, \texttt{--restart=no} evita un
reinicio automático fuera del laboratorio y \texttt{-d} deja el servicio en segundo plano.
El identificador devuelto comenzó con \texttt{59d8457229db}. La figura~\ref{fig:docker}
muestra el contenedor activo y la imagen utilizada.

\labfigure{01_docker_levantamiento.png}
{Contenedor DVWA levantado con la imagen local del laboratorio.}{fig:docker}

\subsection{Redirección de puertos en docker (dvwa)}

Antes de crear el contenedor comprobé que el puerto 6969 no estuviera escuchando. El
mapeo \texttt{127.0.0.1:6969:80} publica el puerto 80 interno de DVWA como 6969 en el
equipo anfitrión. La dirección \texttt{127.0.0.1} es relevante porque limita el acceso a
loopback y evita exponer esta aplicación vulnerable en otras interfaces de red.

Después del primer acceso creé la base de datos desde Setup DVWA, ingresé con el
usuario \texttt{admin} y la contraseña \texttt{password}. Mantuve el nivel de seguridad en
low. En la
figura~\ref{fig:puerto} se ve el servicio respondiendo en la URL local elegida.

\labfigure{02_dvwa_puerto_local.png}
{DVWA disponible en \texttt{127.0.0.1:6969} después de la redirección.}{fig:puerto}

\subsection{Obtención de consulta a replicar (burp)}

Abrí el navegador integrado de burp, inicié sesión en DVWA y activé el interceptor antes
de enviar el formulario Brute Force. La consulta obtenida fue un \texttt{GET} con
los tres parámetros visibles en la URL:

\noindent\begin{minipage}{\textwidth}
\begin{lstlisting}
GET /vulnerabilities/brute/?username=admin&password=password&Login=Login HTTP/1.1
Host: 127.0.0.1:6969
Cookie: PHPSESSID=qiplsqh502ko6d0qlrc7n3jk23; security=low
\end{lstlisting}
\end{minipage}

La cookie conserva la sesión autenticada y el nivel de seguridad bajo. Sin ella, el
formulario protegido no representa la misma prueba. La figura~\ref{fig:burp-raw} muestra
la solicitud completa antes de enviarla a intruder.

\labfigure{06b_burp_consulta_raw.png}
{Consulta del formulario interceptada en burp suite.}{fig:burp-raw}

\subsection{Identificación de campos a modificar (burp)}

En intruder eliminé las posiciones automáticas. El primer campo variable fue
\texttt{username} y el segundo fue \texttt{password}. Dejé el resto de la consulta sin
cambios. Seleccioné cluster bomb porque se necesitaba combinar cada usuario con
cada contraseña, no avanzar ambas listas de forma paralela.

La figura~\ref{fig:burp-posiciones} muestra las dos posiciones delimitadas por
\texttt{\S\ldots\S}. Esta selección produce el producto cartesiano de los dos payloads.

\labfigure{07_burp_posiciones_username_password.png}
{Posiciones de usuario y contraseña configuradas en intruder.}{fig:burp-posiciones}

\subsection{Obtención de diccionarios para el ataque (burp)}

Los usuarios se obtuvieron del índice local \texttt{/hackable/users/}. Los nombres de los
archivos permitieron formar una lista con \texttt{1337}, \texttt{admin}, \texttt{gordonb},
\texttt{pablo} y \texttt{smithy}. Para las contraseñas usé solamente las primeras diez
entradas del archivo rockyou disponible en el equipo.

\begin{center}
\begin{tabular}{ll}
\texttt{123456} & \texttt{princess} \\
\texttt{12345} & \texttt{1234567} \\
\texttt{123456789} & \texttt{rockyou} \\
\texttt{password} & \texttt{12345678} \\
\texttt{iloveyou} & \texttt{abc123}
\end{tabular}
\end{center}

No utilicé el archivo rockyou completo. El subconjunto dejaba una prueba reproducible de
5 usuarios por 10 contraseñas, es decir, 50 combinaciones. Las
figuras~\ref{fig:usuarios} y~\ref{fig:passwords} muestran la fuente de usuarios y el
archivo acotado de contraseñas.

\labfigure[0.62\textwidth]{04_diccionario_usuarios_dvwa.png}
{Usuarios obtenidos desde el índice local de DVWA.}{fig:usuarios}

\labfigure[0.72\textwidth]{05_diccionario_passwords.png}
{Subconjunto de diez contraseñas derivado de rockyou.}{fig:passwords}

\subsection{Obtención de al menos 2 pares (burp)}

La ejecución en intruder cubrió las 50 combinaciones. Ordené los resultados por length porque
el estado HTTP no permitía separarlos, las respuestas correctas e incorrectas devolvieron
\texttt{200}. En el extremo superior aparecieron longitudes distintas a las respuestas
comunes de 4702 o 4703 bytes. La figura~\ref{fig:burp-longitudes} muestra esos candidatos.

\labfigure{10b_burp_resultados_length_desc.png}
{Resultados de intruder ordenados de mayor a menor por longitud.}{fig:burp-longitudes}

No consideré un par válido solo por su longitud. Abrí cada respuesta y comprobé el
mensaje. \texttt{gordonb/abc123} produjo Welcome to the password protected area
gordonb y una longitud de 4745. \texttt{smithy/password} produjo el mismo tipo de
bienvenida para \texttt{smithy} y una longitud de 4742. Las
figuras~\ref{fig:burp-gordon} y~\ref{fig:burp-smithy} muestran ambos resultados.

\labfigure{11_burp_par_gordonb_abc123.png}
{Primer par confirmado en burp, \texttt{gordonb/abc123}.}{fig:burp-gordon}

\labfigure{12_burp_par_smithy_password.png}
{Segundo par confirmado en burp, \texttt{smithy/password}.}{fig:burp-smithy}

Como control, \texttt{gordonb/123456} devolvió Username and/or password
incorrect., sin imagen de usuario y con longitud 4703. La
figura~\ref{fig:burp-invalida} confirma que un estado \texttt{200} no implica una
autenticación correcta.

\labfigure{13_burp_respuesta_invalida.png}
{Respuesta inválida usada para contrastar los pares encontrados.}{fig:burp-invalida}

\subsection{Obtención de código de inspect element (curl)}

Con las herramientas de desarrollador abiertas en la pestaña network, recargué
la página y filtré por documentos. Esto fue necesario porque al principio la tabla estaba
vacía y, en el primer intento de selección, había abierto \texttt{admin.jpg} en vez de la
solicitud del formulario. La comprobación de request URL permitió corregirlo sin
suponer que ambas filas eran equivalentes.

La solicitud correcta aparece en la figura~\ref{fig:network-form}. Sobre esa fila utilicé
copy as curl (bash) y guardé el resultado original en
\texttt{evidencias/raw/curl\_desde\_inspector.sh}. La figura~\ref{fig:curl-codigo} muestra
el inicio del comando, incluida la URL, la cookie y el \texttt{Referer}.

\labfigure{14b_network_solicitud_formulario.png}
{Solicitud principal seleccionada en network antes de copiarla como curl.}{fig:network-form}

\labfigure{25b_curl_codigo_inspect.png}
{Inicio del código curl obtenido desde las herramientas de desarrollador.}{fig:curl-codigo}

\subsection{Utilización de curl por terminal (curl)}

Conservé el comando copiado sin editar y preparé una segunda copia donde cambié solamente
\texttt{password=password} por \texttt{password=clave\_incorrecta\_lab2}. Las dos
ejecuciones fueron:

\begin{lstlisting}
bash evidencias/raw/curl_desde_inspector.sh > evidencias/raw/curl_valido.html
bash evidencias/raw/curl_invalido.sh > evidencias/raw/curl_invalido.html
\end{lstlisting}

El primer archivo contiene Welcome to the password protected area admin. El
segundo contiene Username and/or password incorrect. La
figura~\ref{fig:curl-ejecuciones} muestra los dos comandos registrados por la terminal. El
tráfico quedó guardado aparte en \texttt{curl\_valido\_invalido.pcapng}.

\labfigure{26_curl_ejecuciones_terminal.png}
{Ejecuciones de curl válida e inválida realizadas desde la terminal.}{fig:curl-ejecuciones}

\subsection{Demuestra 5 diferencias (curl)}

La comparación se realizó sobre archivos obtenidos realmente, no sobre páginas
reconstruidas. Encontré cinco diferencias.

\begin{enumerate}
  \item En la solicitud válida, el campo \texttt{password} contiene \texttt{password}.
  En la inválida contiene \path{clave_incorrecta_lab2}.
  \item El cuerpo válido mide 4413 bytes y el inválido 4375, una diferencia de 38 bytes.
  \item La respuesta válida muestra una bienvenida para \texttt{admin}; la inválida
  muestra el mensaje de credenciales incorrectas.
  \item La válida presenta el resultado dentro de \texttt{<p>}; la inválida usa
  \texttt{<pre><br />}. 
  \item La válida incluye \texttt{/hackable/users/admin.jpg}; la inválida no contiene
  una imagen de usuario.
\end{enumerate}

La figura~\ref{fig:curl-entrada} demuestra que el parámetro de contraseña fue la única
modificación en la solicitud. La figura~\ref{fig:curl-tamanos} conserva los tamaños y la
figura~\ref{fig:curl-diff} muestra la línea HTML distinta.

\labfigure{18_curl_diferencia_solicitud.png}
{Única modificación entre la solicitud válida y la inválida.}{fig:curl-entrada}

\labfigure{16_curl_tamanos_respuestas.png}
{Tamaños de los cuerpos guardados para ambas respuestas.}{fig:curl-tamanos}

\labfigure{17_curl_diferencias_contenido.png}
{Diferencia de contenido entre la respuesta inválida y la válida.}{fig:curl-diff}

En ambos casos el servidor respondió \texttt{200}. Esto no cuenta como diferencia, pero
sí muestra por qué el resultado debe decidirse con el cuerpo y no solo con el código de
estado.

\subsection{Instalación y versión a utilizar (hydra)}

La herramienta hydra no estaba instalada al comenzar. El comando \texttt{command -v hydra} no devolvió
una ruta, mientras que el índice local de Debian sí mostraba el paquete. Lo instalé con:

\begin{lstlisting}
sudo apt install hydra
\end{lstlisting}

Después confirmé \texttt{/usr/bin/hydra} y la versión 9.4. La
figura~\ref{fig:hydra-version} conserva esa verificación.

\labfigure{24_hydra_version.png}
{Versión de hydra utilizada en el laboratorio.}{fig:hydra-version}

\subsection{Explicación de comando a utilizar (hydra)}

Antes de ejecutar revisé la ayuda local de \texttt{http-get-form}. El módulo exige la ruta,
los parámetros y una condición para reconocer el fallo. Usé el siguiente comando.

\begin{lstlisting}
hydra -L evidencias/diccionarios/users_dvwa_5.txt \
  -P evidencias/diccionarios/passwords_rockyou_10.txt \
  127.0.0.1 -s 6969 http-get-form \
  "/vulnerabilities/brute/:username=^USER^&password=^PASS^&Login=Login:H=Cookie\: security=low; PHPSESSID=qiplsqh502ko6d0qlrc7n3jk23:F=Username and/or password incorrect" -V
\end{lstlisting}

\begin{itemize}
  \item \texttt{-L} y \texttt{-P} cargan las listas de cinco usuarios y diez
  contraseñas.
  \item \texttt{127.0.0.1 -s 6969} fija el único destino y puerto autorizados.
  \item \texttt{http-get-form} corresponde al método observado en burp y network.
  \item \verb|^USER^| y \verb|^PASS^| son reemplazados por cada combinación.
  \item \texttt{H=Cookie} mantiene la sesión y \texttt{security=low}.
  \item \texttt{F=Username and/or password incorrect} define el texto de una respuesta
  fallida.
  \item \texttt{-V} muestra cada intento. No se agregó \texttt{-f}, porque detenerse en
  el primer hallazgo impediría cumplir la búsqueda de al menos dos pares.
\end{itemize}

El máximo seguía siendo 50 combinaciones. No se usó el rockyou completo ni se apuntó a
otra dirección.

\subsection{Obtención de al menos 2 pares (hydra)}

La ejecución de hydra inició a las 23:42:54 y terminó a las 23:42:55. Informó tres pares válidos.

\begin{itemize}
  \item \texttt{admin/password}
  \item \texttt{gordonb/abc123}
  \item \texttt{smithy/password}
\end{itemize}

La figura~\ref{fig:hydra-resultados} contiene las líneas originales y el resumen
\texttt{3 valid passwords found}. Los dos últimos pares coinciden con los confirmados en
burp. Que 50 intentos pudieran completarse sin bloqueo también muestra la ausencia de un
límite efectivo en esta configuración de DVWA.

\labfigure{19_hydra_resultados.png}
{Tres pares válidos encontrados por hydra.}{fig:hydra-resultados}

\subsection{Explicación paquete curl (tráfico)}

Seleccioné la trama 978 porque corresponde a la solicitud válida y está enlazada con la
respuesta 980. En la figura~\ref{fig:paquete-curl} se observa un \texttt{GET} mediante
HTTP/1.1, desde el puerto origen 51784 hacia el 6969. La trama mide 871 bytes e incluye la
URL completa, la cookie y varios encabezados del navegador.

Aunque la solicitud fue emitida por curl, el \texttt{User-Agent} indica chromium 150. No
es un error de wireshark, Copy as curl conservó explícitamente ese encabezado. La
respuesta 980 tiene estado 200, \texttt{Content-Length: 4413}, no declara compresión y mide
4751 bytes como trama capturada.

\labfigure{21_curl_paquete_http.png}
{Detalle HTTP de la solicitud válida emitida mediante curl.}{fig:paquete-curl}

\subsection{Explicación paquete burp (tráfico)}

Para burp seleccioné la trama 402, que contiene el par válido
\texttt{gordonb/abc123}. La solicitud mide 875 bytes, usa \texttt{GET} mediante HTTP/1.1
y conserva el mismo User-Agent de chromium. A diferencia de la captura curl, aquí aparecen
\texttt{Accept-Encoding: gzip, deflate, br} y \texttt{Connection: keep-alive}. La
figura~\ref{fig:paquete-burp} muestra esos campos y el vínculo con la trama de respuesta
404.

\labfigure{22_burp_paquete_http.png}
{Solicitud válida enviada por burp intruder.}{fig:paquete-burp}

La trama de respuesta 404, mostrada en la figura~\ref{fig:respuesta-burp}, tiene estado
HTTP 200 y usa compresión \texttt{gzip}. Su \texttt{Content-Length} es 1476. Según wireshark,
ese cuerpo se convierte en 4417 bytes al decodificarlo. Por eso no comparé directamente
ese valor con la longitud 4745 mostrada por intruder, que corresponde a otra
representación de la respuesta.

\labfigure{23_burp_respuesta_gzip.png}
{Respuesta comprimida asociada a la solicitud de burp.}{fig:respuesta-burp}

\subsection{Explicación paquete hydra (tráfico)}

La trama 2338 corresponde al intento con usuario \texttt{admin} y contraseña
\texttt{password}. Es bastante menor que las anteriores, 267 bytes. La herramienta hydra usa \texttt{GET}
mediante HTTP/1.0 y envía \texttt{User-Agent: Mozilla/5.0 (Hydra)}. En la cookie aparece
primero \texttt{security=low} y después \texttt{PHPSESSID}. No aparece
\texttt{Accept-Encoding}. La trama de respuesta 2394 tiene estado 200, cuerpo de 4413 bytes,
\texttt{Connection: close} y no declara compresión.

La figura~\ref{fig:paquete-hydra} permite reconocer la herramienta directamente por el
User-Agent y la versión HTTP. Además, la captura completa contiene 100 solicitudes, 50
consultas previas a la ruta sin parámetros y 50 intentos con credenciales, cada una en un
flujo TCP distinto.

\labfigure{20_hydra_paquete_http.png}
{Detalle HTTP de un intento válido emitido por hydra.}{fig:paquete-hydra}

\subsection{Mención de las diferencias (tráfico)}

La tabla~\ref{tab:trafico} resume campos extraídos de los tres PCAPNG. Los conteos no son
una prueba de rendimiento aislada, porque cada captura representa una actividad distinta,
pero sí describen el comportamiento observado en este laboratorio.

\begin{table}[htbp]
\centering
\label{tab:trafico}
\small
\begin{tabularx}{\textwidth}{|>{\raggedright\arraybackslash}p{3.1cm}|X|X|X|}
\hline
\textbf{Característica} & \textbf{curl} & \textbf{burp} & \textbf{hydra} \\
\hline
Trama analizada & 978 & 402 & 2338 \\
\hline
Método y versión & GET, HTTP/1.1 & GET, HTTP/1.1 & GET, HTTP/1.0 \\
\hline
User-Agent & chromium 150 & chromium 150 & Mozilla/5.0 (Hydra) \\
\hline
Tamaño de solicitud & 871 bytes & 875 bytes & 267 bytes \\
\hline
Codificación aceptada & No detectada & gzip, deflate, br & No detectada \\
\hline
Respuesta & 4413 bytes, sin compresión & 1476 bytes transferidos con gzip & 4413 bytes, sin compresión \\
\hline
Conexión observada & Dos solicitudes en dos flujos & 53 solicitudes en 38 flujos & 100 solicitudes en 100 flujos \\
\hline
Patrón de trabajo & Una ejecución válida y una inválida & 50 combinaciones y tres repeticiones & 50 consultas previas y 50 intentos \\
\hline
\end{tabularx}
\caption{Comparación del tráfico observado.}
\end{table}

Las tres herramientas enviaron credenciales dentro de la URL porque el formulario usa
GET. Las dos pruebas solicitadas se realizaron con curl. La herramienta burp mantuvo conexiones y solicitó
respuestas comprimidas. La herramienta hydra abrió más flujos, generó una consulta previa por cada intento
y envió solicitudes mucho más pequeñas. En burp se contaron 53 solicitudes con parámetros,
las 50 combinaciones, una aparición adicional de \texttt{admin/password} y dos apariciones
adicionales de \texttt{gordonb/iloveyou}. La captura demuestra la repetición, pero no permite
atribuirle una causa con seguridad.

\subsection{Detección de SW (tráfico)}

La herramienta hydra es la más fácil de detectar en estas capturas. Su User-Agent contiene el
nombre \texttt{Hydra}, usa HTTP/1.0, cierra la conexión y genera una secuencia de consultas
previas e intentos en flujos separados. No hace falta inferir su origen solo por el tamaño.

La situación cambia con curl y burp. Ambos paquetes anuncian chromium 150 porque las dos
herramientas conservaron encabezados del navegador. Por lo tanto, sería incorrecto afirmar
que el User-Agent permite reconocerlas. Dentro de este experimento se pueden separar por la
compresión, los encabezados y el patrón de conexiones, burp pidió gzip y reutilizó parte de
los flujos, mientras el curl guardado no pidió compresión y solo realizó dos solicitudes.
Esas diferencias describen estas ejecuciones, pero no son una firma universal. Un comando
Curl puede cambiar sus encabezados y burp también puede modificarlos.

\section*{Conclusiones y comentarios}
\addcontentsline{toc}{section}{Conclusiones y comentarios}

DVWA quedó aislado en \texttt{127.0.0.1:6969} y las tres herramientas trabajaron contra
el mismo formulario. En burp se confirmaron \texttt{gordonb/abc123} y
\texttt{smithy/password}. La ejecución con hydra encontró esos dos pares y también
\texttt{admin/password} después de probar las 50 combinaciones. En curl, la respuesta
válida midió 4413 bytes y la inválida 4375. Aun así, ambas devolvieron HTTP 200. Ese fue el
resultado que más me llamó la atención, porque obliga a revisar el mensaje y la estructura
del cuerpo en vez de confiar en el estado.

El problema más evidente es que el formulario usa GET. Usuario y contraseña quedan en la
URL, por lo que aparecen en burp, en el comando curl y en wireshark. También fue posible
probar todo el diccionario sin bloqueo ni espera. Con solo cinco usuarios y diez
contraseñas aparecieron tres pares válidos. En una aplicación real, esa combinación de
credenciales débiles y ausencia de límite de intentos facilita un ataque de diccionario muy
pequeño.

La herramienta burp fue la que preferí porque ya la había usado en CTF y su vista de
solicitudes y respuestas me resultó familiar. La herramienta hydra fue mucho más directa para automatizar
la prueba. La advertencia que muestra al iniciar sobre usos militares fue una de las pocas
cosas del procedimiento que me causó gracia. En el tráfico también fue la herramienta más
fácil de reconocer, mientras curl y burp conservaron el mismo User-Agent de chromium.

La parte más difícil fue seguir la actividad. La explicación presencial del profesor fue
poco clara y el video resultó todavía peor como guía, así que terminé buscando la sintaxis y
comprobando cada paso por mi cuenta. Al igual que en el laboratorio anterior, no encontré
un aprendizaje técnico nuevo. La actividad fue mecánica y aburrida. Lo que sí exigió
cuidado fue no aceptar resultados por apariencia, ordené burp por longitud, comprobé el
texto de cada respuesta y vinculé cada paquete con su trama de respuesta antes de escribir
las conclusiones.

\section*{Anexo}
\addcontentsline{toc}{section}{Anexo}

\begin{itemize}
  \item \href{https://github.com/lavapatos/cripto-labs/tree/main/lab-2}{Repositorio}
\end{itemize}

\end{document}
